% !PRINTED 451
\weberpage{451}
%<<CONT
of the second equation reduces to $D$ only when it is imprimitive and
\weq{8}{\begin{gathered}
B' - 2A'c \ \text{by} \ p, \\
A'^{2}c - B'c + C' \ \text{by} \ p^2 % [sic]
\end{gathered}}
is divisible, and then the corresponding number (6) satisfies at the same time equation (5) and equation (2). But this is the case, as we shall show at once, only for one value $c$, and consequently the value $u$ in question can be expressed \gesperrt{rationally} in terms of $v$.

For if in the first place $p$ is odd, the first congruence (8):
\[
2A'c - B' \equiv 0 \pmod{p}
\]
has only \gesperrt{one} root, and for this
\[
(2A'c - B')^2 - p^2D = 4A'(A'c^2 - B'c + C') \equiv 0 \pmod{p^2}.
\]
If $p = 2$, then $D'$, and consequently $B'$, is even, and hence $B' - 2A'c$ is always divisible by 2. We then have
\[
\left(A'c - \frac{B'}{2}\right)^2 - D = A'(A'c^2 - B'c + C'),
\]
and since $D \equiv 0$ or $\equiv 1 \pmod{4}$, the left side here is divisible by 4 if it is assumed that $A'c - \frac{B'}{2} \equiv D \pmod{2}$.

If, conversely, one starts from an equation
\weq{9}{A\omega^2 + B\omega + C = 0, \quad B^2 - 4AC = D}
of discriminant $D$, takes $A$ relatively prime to $p$, and sets $\omega' = p\omega$, then $\omega'$ satisfies the primitive equation
\weq{10}{A\omega'^2 + Bp\omega' + Cp^2 = 0}
of discriminant $D' = p^2D$; and if we therefore set $v = j(\omega')$
then $u = j(\omega)$ is rationally expressible in terms of $v$. From this, then, the theorem follows:

1. \gesperrt{Every class invariant for the discriminant $D$ is rationally expressible in terms of a class invariant for the discriminant $p^2D$.}

In order, however, to answer the question how many values of $v$ lead to the same value of $u = j(\omega)$, we observe that the various values of $v$ which do this must be roots of equation (5) and (3), and hence must have one of the forms
\weq{11}{j(p\omega), \quad j\left(\frac{\omega + c}{p}\right)}
%CONT>>
